\section{Proofs for Section~\ref{sec:optsets}}\label{app:proximal}

This appendix contains the proofs for Section~\ref{sec:optsets}, the lemma on
proximal stages used by DISC (Algorithm~\ref{alg:disc}), and
Proposition~\ref{prop:sshard} on the hardness of discovery within
$\mathfrak D$.

\subsection{One center and two minimizers}

\begin{proof}[Proof of Proposition~\ref{prop:twocenters}]
\emph{The instance.} $F_M=x^2+z(M-2x)$. For $x<M/2$ the minimum over $z$ is
at $z=0$ with value $x^2$; for $x>M/2$ it is at $z=M$ with value $(M-x)^2$;
at $x=M/2$ the value is $M^2/4$. So $\mathcal S$ and $\OPT$ are as stated,
and the Hessian $\bigl(\begin{smallmatrix}2&-2\\-2&0\end{smallmatrix}\bigr)$
gives $L_x=2$, $L_z=0$. The map $(x,z)\mapsto(M-x,M-z)$ preserves $F_M$ (a
direct expansion) and swaps the minimizers, so for growth we may assume
$x\le M/2$. If $x\le M/4$, then $M-2x\ge M/2\ge z/2$ and
$F_M\ge x^2+z^2/2\ge\frac12\norm{(x,z)}^2$. If $M/4<x\le M/2$, then
$F_M\ge x^2\ge M^2/16$ while $\norm{(x,z)}^2\le5M^2/4$. In both cases
$F_M\ge\frac1{20}\dist((x,z),\mathcal S)^2$, and
$\kappa_S\le L/g_S=40$ with $g_S=1/20$.

\emph{The bound~\eqref{eq:twocenters}.} By symmetry assume $M-c_x\ge M/2$
(otherwise use the reflected argument with $z=0$ and the endpoint $0$). Let
$c_x=v_0<v_1<\dots<v_k=M$ be the $x$-nodes on the right of the center, so
$k\ge1$, and put $\delta_m=v_m-c_x$. All steps except the last are unclipped:
$\delta_{m+1}=(1+\theta)\delta_m+h$ for $m<k-1$, and
$\delta_k\le(1+\theta)\delta_{k-1}+h$. The node $M$ lies in $G_z$,
$F_M(v,M)=(M-v)^2$, every correction is nonnegative, and
$d_x(v)\ge w_x(v)^2/4$.

If $k=1$, the node $M$ has the adjacent interval $[c_x,M]$ of length at least
$M/2$, so $\beta\le Q(M,M)\le-M^2/16$, and
$M^2/16\ge\max\{\theta^2M^2/379,h^2/20\}$ since $\theta\le1/4$ and $h\le M$.

If $k\ge2$, put $\alpha=v_{k-1}-v_{k-2}$ and $\alpha'=M-v_{k-1}$. Then
$Q(M,M)\le-\alpha'^2/4$ and $Q(v_{k-1},M)\le\alpha'^2-\alpha^2/4$. If
$\alpha'^2\ge\alpha^2/5$, the first is at most $-\alpha^2/20$; otherwise the
second is below $-\alpha^2/20$; so $\beta\le-\alpha^2/20$. Now
$\alpha=h+\theta\delta_{k-2}\ge h$, and the recursion gives
$\delta_{k-2}\ge(\delta_k-(2+\theta)h)/(1+\theta)^2$. If $h\le M/16$, then
with $\delta_k\ge M/2$ and $\theta\le1/4$,
$\delta_{k-2}\ge(\frac M2-\frac{9M}{64})\cdot\frac{16}{25}=\frac{23M}{100}$, so
$\alpha\ge\frac{23}{100}\theta M$. If $h>M/16$, then
$\alpha\ge h>M/16\ge\theta M/4$. Hence $\alpha\ge\max\{h,\frac{23}{100}\theta M\}$
and $\beta\le-\max\{h^2,(0.23\,\theta M)^2\}/20$, which
implies~\eqref{eq:twocenters} because $0.23^2/20>1/379$.

\emph{Consequences.} (a) Since $L_z=0$, coordinate $z$ is not in $P$: TRIAL
uses $G_z=\{0,M\}$ and never filters $z$. Filtering with $U\ge\OPT$ keeps
both minimizers (Proposition~\ref{prop:filter}), so the $x$-interval of every
stage is $[0,M]$. At a stage $j\ge1$ the mesh is $h_{xj}=h_{x0}2^{-j}<M$ and
the center lies in $X$, so~\eqref{eq:twocenters} applies. At stage $0$ the
center is $\ell=(0,0)$ and $M\le h_{x0}<2M$ (Section~\ref{sec:growth}), so
$G_x=\{0,M\}$, $d_x(M)\ge M^2/4$ and $\beta\le Q(M,M)\le-M^2/4$; this is at
most $-\max\{\theta^2M^2/379,h_{x0}^2/20\}$ because $h_{x0}^2/20<M^2/5$.
With a common mesh, $h_j=M2^{-j}\le M$ at every stage.

(b) By~(a), $\beta\ge-\varepsilon$ gives $\theta M\le\sqrt{379\varepsilon}$
and $h\le\sqrt{20\varepsilon}$, so every step of the $x$-grid is at most
$h+\theta M\le(\sqrt{20}+\sqrt{379})\sqrt\varepsilon<24\sqrt\varepsilon$. The
side of $[0,M]$ from $c_x$ to the farther endpoint has length at least $M/2$,
so it takes at least $M/(48\sqrt\varepsilon)$ steps and contains more than
$M/(48\sqrt\varepsilon)$ nodes. Since $\abs{G_z}=2$, the table of the bag
$\{x,z\}$ has more than $M/(24\sqrt\varepsilon)$ entries.

(c) Apply (b) with $\varepsilon=1$. EX stops only when
$F(\tilde x)-\beta<1/(\Omega W)\le1$ for the bound $\beta$ of the successful
stage of CT, and $F(\tilde x)\ge\OPT=0$, so that stage has $\beta>-1$. The
data of $F_M$ are $1$, $-2$, $M$ and the endpoints $0$ and $M$, so
$I=O(\log M)$.
\end{proof}

\subsection{Coordinatewise concave quadratics}

\begin{proof}[Proof of Lemma~\ref{lem:endpointid}]
Since $B_{t'}^{\uparrow}\subseteq B_t$ for a child $t'$, the bracket is a
function of $v_{B_t}$, and $r_t\ge0$ because $M_t$ is its minimum over the
labels of $B_t\setminus B_t^\uparrow$.

\emph{Telescoping.} Let $v\in E([n])$. Every nonroot table $M_t$ occurs in
$\sum_tr_t(v_{B_t})$ once with a plus sign, in the residual of its parent, and
once with a minus sign. Hence $\sum_tr_t(v_{B_t})=\sum_tq_t(v_{B_t})-M_{t_0}
=F(v)-M_{t_0}$.

\emph{Expectation.} Fix $x\in\bar X$ and let $Y$ be the endpoint rounding
of Section~\ref{sec:endpointset}. Then $\E Y_i=x_i$,
$\Var Y_i=(x_i-\ell_i)(u_i-x_i)$, $\E Y_iY_j=x_ix_j$ for $i\ne j$,
and $\E Y_i^2=x_i^2+\Var Y_i$. Therefore
$\E F(Y)=F(x)+\frac12\sum_iH_{ii}(x_i-\ell_i)(u_i-x_i)$, while
$\E r_t(Y_{B_t})=\sum_vr_t(v)\pi_t(v;x)$. Taking expectations in the
telescoping identity gives \eqref{eq:endpointid} with $M_{t_0}$ in place
of $\OPT$.

\emph{Value.} Every term on the right of \eqref{eq:endpointid} is
nonnegative on $\bar X$, so $F\ge M_{t_0}$ on $\bar X\supseteq X$. Trace back:
choose labels of $B_{t_0}$ attaining $M_{t_0}$; for each child $t'$ of a
processed bag $t$, choose the labels of $B_{t'}\setminus B_{t'}^\uparrow$
attaining $M_{t'}(v_{B_{t'}^\uparrow})$. By running intersection, a variable
belongs to $B_t\setminus B_t^\uparrow$ exactly for the bag $t$ nearest the
root among those containing it, and to $B_{t}^{\uparrow}$ for every other bag
containing it. Hence every variable receives exactly one label, and the
resulting endpoint assignment $v^*$ has all residuals zero. By telescoping
$F(v^*)=M_{t_0}$. Both endpoints of every $X_i$ are feasible, so $v^*\in X$
and $\OPT=M_{t_0}$.

\emph{Checking.} The argument uses only the nonnegativity of the residuals
and one endpoint assignment with zero residuals, so any tables with these two
properties give \eqref{eq:endpointid}.

\emph{Cost.} A bag table has at most $2^p$ entries, and children are combined
by addition. By induction on the tree, every message entry is the sum of the
assigned factors over the subtree at one endpoint assignment. Its bit length
is therefore bounded by the bit length of a common denominator of all endpoint
monomial values plus $\log_2$ of the number of monomials times their maximum
magnitude, both polynomial in $I$.
\end{proof}

\begin{proof}[Proof of Corollary~\ref{cor:facecsp}]
(i) Condition (a) of Theorem~\ref{thm:endpointset} says $\chi_i(x)\ne\ast$
when $H_{ii}<0$; and $\chi_i(x)=\ast$ is possible only if $X_i$ has a point
strictly between its endpoints. Condition (b) says that every $v$ with
$\pi_t(v;x)>0$ has $r_t(v)=0$, that is,
$\operatorname{cube}(\chi(x)_{B_t})\subseteq Z_t$.

(ii) Replacing $\ast$ by an endpoint shrinks every cube. If $\chi$ is
admissible and $x\in X\cap\operatorname{face}(\chi)$, then $\chi(x)$ is
obtained from $\chi$ by replacing some $\ast$ entries by endpoints, so $x$ is
optimal by (i). Conversely $x\in\mathcal S$ lies in
$\operatorname{face}(\chi(x))$. A pattern $\chi$ determines the set
$\{x\in X:\chi(x)=\chi\}$, which is nonempty by the definition of $\Sigma_i$.

(iii) Each $\mathrm{Adm}_t$ is computed by testing at most $3^p$ patterns
against at most $2^p$ residual entries. A constraint problem whose relations
live on the bags of a tree decomposition is solved by nonserial dynamic
programming over $(T,(B_t))$ in the min-sum, max-sum and sum-product
semirings, with $O(p\,3^p)$ operations per bag \cite{Dechter1999}; each unary
cost is charged once, at the bag nearest the root that contains its variable.
The minimizer is unique if and only if exactly one pattern is admissible and
it has no $\ast$, by the disjointness in (ii). The dimension of $\mathcal S$
is the maximum number of $\ast$ entries on continuous coordinates. When
$\mathcal S$ is finite,
$|\mathcal S|=\sum_\chi\prod_{i:\chi_i=\ast}|X_i\cap(\ell_i,u_i)|$. By (ii),
$\dist(y,\mathcal S)^2$ is the minimum over admissible $\chi$ of the sum of
the unary costs $(y_i-\ell_i)^2$, $(y_i-u_i)^2$ and $\dist(y_i,X_i)^2$ for
$\chi_i=\mathsf L$, $\mathsf U$ and $\ast$. A linear objective attains its
minimum over each $X\cap\operatorname{face}(\chi)$ at an endpoint pattern, so
it is minimized by the same recursion over $\{\mathsf L,\mathsf U\}$ labels.
All numbers are rationals of polynomial bit length.
\end{proof}

\subsection{The diagonal certificate and the proximal stage}

As in Section~\ref{sec:diagcert}, $X$ is a continuous box and $F$ a rational
quadratic.

\begin{proof}[Proof of Lemma~\ref{lem:diagcert}]
(i) Both sides are quadratic polynomials in $x$ that vanish at $\bar x$.
Their Hessians are $H$ and $H+\Lambda-\Lambda=H$. The $i$th partial
derivative of the right side at $\bar x$ is
$\frac12\lambda_i(\bar x)(u_i+\ell_i-2\bar x_i)$. This equals
$|\partial_iF(\bar x)|=\partial_iF(\bar x)$ when $\bar x_i=\ell_i$, equals
$-|\partial_iF(\bar x)|=\partial_iF(\bar x)$ when $\bar x_i=u_i$, and is
$0=\partial_iF(\bar x)$ at an interior coordinate. Two quadratics with equal
value, gradient and Hessian at one point coincide.

(ii) On $X$ both terms on the right of \eqref{eq:diagid} are nonnegative,
so $\bar x$ is optimal, and $x\in X$ is optimal exactly when both vanish. A
positive semidefinite quadratic form vanishes exactly on the kernel of its
matrix.

(iii) For $x\in X$ the bracket in $\Psi$ is at most $F(x)$, so
$\Psi(\lambda)\le\OPT$. If (a) holds at $\bar x$, then by (i) the bracket
with $\lambda=\lambda(\bar x)$ equals
$F(\bar x)+\frac12(x-\bar x)^{\T}(H+\Lambda(\bar x))(x-\bar x)$, whose
infimum over $\R^n$ is $F(\bar x)=\OPT$; this gives (c). Suppose (c) holds
for some $\lambda$, and let $\bar x\in\mathcal S$ be arbitrary. Write
$\mathcal L(x)$ for the bracket. Then
$\OPT=\Psi(\lambda)\le\mathcal L(\bar x)\le F(\bar x)=\OPT$. Hence
$\lambda_i(\bar x_i-\ell_i)(u_i-\bar x_i)=0$ for all $i$, and $\bar x$
minimizes the quadratic $\mathcal L$ over $\R^n$. So
$\nabla^2\mathcal L=H+\diag(\lambda)\succeq0$ and
$0=\partial_i\mathcal L(\bar x)=\partial_iF(\bar x)+\lambda_i\bigl(\bar x_i-\tfrac{\ell_i+u_i}2\bigr)$.
At an interior coordinate, complementarity gives
$\lambda_i=0=\lambda_i(\bar x)$. At $\bar x_i=\ell_i$ the stationarity
equation gives $\lambda_i=2\partial_iF(\bar x)/(u_i-\ell_i)=\lambda_i(\bar x)$,
and similarly at $\bar x_i=u_i$. Thus $\lambda=\lambda(\bar x)$ and
$H+\Lambda(\bar x)\succeq0$. Since $\bar x\in\mathcal S$ was arbitrary, this
proves (b) and the uniqueness and invariance statements. Finally (b) implies
(a) because $\mathcal S\ne\emptyset$.
\end{proof}

\begin{example}[Tilted and disconnected optimal sets in $\mathfrak D$]\label{ex:diagtilt}
Let $F=(x-y+t/2)^2+\frac18t(1-t)$ on $[0,1]^3$, and $a=(1,-1,\frac12)$. Both
terms are nonnegative on the box and vanish together exactly when
$t\in\{0,1\}$ and $x-y+t/2=0$. So $\OPT=0$, and $\mathcal S$ consists of the
two segments $\{(v,v,0):0\le v\le1\}$ and $\{(v,v+\frac12,1):0\le
v\le\frac12\}$. The Hessian $H=2aa^{\T}-\diag(0,0,\frac14)$ has diagonal
$(2,2,\frac14)$ and is indefinite, since $w=(0,1,2)$ has $a^{\T}w=0$ and
$w^{\T}Hw=-1$. At every minimizer $\partial_xF=\partial_yF=0$ and
$\partial_tF=\frac18(1-2t)$, which is $\frac18$ at $t=0$ and $-\frac18$ at
$t=1$. So every minimizer is a box KKT point with
$\Lambda=\diag(0,0,\frac14)$, and $H+\Lambda=2aa^{\T}\succeq0$; hence
$F\in\mathfrak D$.
\end{example}

\begin{lemma}[Proximal stages]\label{lem:proximal}
Let $K_\theta=8\theta^{-1}\lceil\log_2(n+2)\rceil$. In PROX$(\hat\kappa)$
(Algorithm~\ref{alg:prox}), $\theta\le1/4$, $8\hat\kappa\theta^2\le1$ and
$\theta^{-1}\le6\sqrt{\hat\kappa}$, and:
\begin{enumerate}[label=(\roman*)]
\item every grid has at most $K_\theta$ nodes per coordinate;
\item if $F$ has set growth with a constant $g_S$ such that
$\kappa_S\le\hat\kappa$, then a proximal stage whose center satisfies
$\dist(c,\mathcal S)^2\le4\hat\kappa nh^2$ returns $y^+$ with
$F(y^+)-\OPT\le\frac{99}{256}Lnh^2$ and
$\dist(y^+,\mathcal S)^2\le\frac{99}{256}\hat\kappa nh^2$;
\item under the hypothesis of (ii) on $g_S$,
$\dist(y^{(j)},\mathcal S)^2\le\frac{99}{256}\hat\kappa nh_j^2$ for every
$j\ge0$;
\item stage $j$ costs $O\bigl(p(\abs{\mathcal A}+n+N)K_\theta^p\bigr)$ table
operations on rationals of bit length $O(I+j+\mu K_\theta)$.
\end{enumerate}
\end{lemma}

\begin{proof}[Proof of Lemma~\ref{lem:proximal}]
\emph{Constants.} Since $\mu\ge2$, $\theta\le1/4$, and
$4^{-\mu}\hat\kappa\le1/8$ is $8\hat\kappa\theta^2\le1$. If $\mu=2$, then
$\theta^{-1}=4\le6\sqrt{\hat\kappa}$. Otherwise minimality of $\mu$ gives
$4^{1-\mu}\hat\kappa>1/8$, so $\theta^{-1}=2^\mu<2\sqrt{8\hat\kappa}
<6\sqrt{\hat\kappa}$.

(i) Each endpoint of $X'_i$ is within $\varrho h$ of $c_i$, so by
Lemma~\ref{lem:graded}(b) each side of $G_i$ has at most
$\lceil(4/\theta)\ln(1+\theta\varrho)\rceil$ nodes besides $c_i$. Now
$\theta\varrho\le2\theta\sqrt{\hat\kappa n}+2\theta\le\sqrt{n/2}+\frac12$,
because $\theta\sqrt{\hat\kappa}\le1/\sqrt8$ and $\theta\le\frac14$, and
$\frac32+\sqrt{n/2}\le n+2$. Hence
$\abs{G_i}\le1+2\lceil(4/\theta)\ln(n+2)\rceil\le3+(8/\theta)\ln(n+2)$. With
$3\le3/(4\theta)$, $8\ln(n+2)\le5.6\log_2(n+2)$ and
$\lceil\log_2(n+2)\rceil\ge2$, this is at most
$7\theta^{-1}\lceil\log_2(n+2)\rceil\le K_\theta$.

(ii) Let $x^\circ\in\mathcal S$ be nearest to $c$ (it exists since $\mathcal S$
is compact), and put $\delta=\norm{x^\circ-c}$, so
$\delta^2\le4\hat\kappa nh^2$. Each
$|x^\circ_i-c_i|\le\delta\le2\sqrt{\hat\kappa n}\,h\le\varrho h$, so
$x^\circ\in X'$. For each $i$ choose a grid interval $J_i$ of $G_i$ that
contains $x^\circ_i$ and
whose endpoint $a_i$ nearer to $c_i$ satisfies $|a_i-c_i|\le|x^\circ_i-c_i|$:
if $x^\circ_i>c_i$, the interval $[a_i,a_i']$ with $a_i<x^\circ_i\le a_i'$; if
$x^\circ_i<c_i$, symmetrically; if $x^\circ_i=c_i$, any interval with endpoint
$c_i$. By Lemma~\ref{lem:graded}(a) at the node $a_i$,
$w(J_i)\le h+\theta|x^\circ_i-c_i|$. Let $C=\prod_iJ_i$ and
$\tilde F(y)=F(y)+\omega\norm{y-c}^2$, which has the upper coordinate
curvatures $L_i+2\omega$. The proof of Proposition~\ref{prop:cellwise}(a) uses
only Definition~\ref{def:curvature}, so it applies to $\tilde F$ and gives a
vertex $v$ of $C$ with
$\tilde F(v)-\sum_i\frac{L_i+2\omega}8w(J_i)^2\le\tilde F(x^\circ)=\OPT+\omega\delta^2$.
Since $v_i$ is an endpoint of $J_i$, $d_i(v_i)\ge L_iw(J_i)^2/8$, and so
\[
Q_\omega(y^+)\le Q_\omega(v)\le\tilde F(v)-\sum_i\frac{L_i}8w(J_i)^2
\le\OPT+\omega\delta^2+\frac\omega4\sum_iw(J_i)^2 .
\]
With $w(J_i)^2\le2h^2+2\theta^2(x^\circ_i-c_i)^2$ this gives
\[
Q_\omega(y^+)\le\OPT+\omega\Bigl[\bigl(1+\tfrac{\theta^2}2\bigr)\delta^2
+\tfrac{nh^2}2\Bigr].
\]
By Lemma~\ref{lem:graded}(a), $w_i(v)\le h+\theta|v-c_i|$ for every node $v$,
so $d_i(v)\le\frac L4h^2+\omega|v-c_i|^2$ and
$D(y^+)\le\frac L4nh^2+\omega\norm{y^+-c}^2$. Therefore
\begin{align*}
F(y^+)-\OPT&=Q_\omega(y^+)-\OPT+D(y^+)-\omega\norm{y^+-c}^2\\
&\le\frac{Lnh^2}4\Bigl[1+4\hat\kappa\theta^2\bigl(1+\tfrac{\theta^2}2\bigr)
+\tfrac{\theta^2}2\Bigr]\le\frac{Lnh^2}4\cdot\frac{99}{64},
\end{align*}
using $\delta^2\le4\hat\kappa nh^2$, $\omega=L\theta^2/4$,
$4\hat\kappa\theta^2\le\frac12$ and $\theta^2\le\frac1{16}$. Set growth gives
$\dist(y^+,\mathcal S)^2\le(F(y^+)-\OPT)/g_S\le\frac{99}{256}(L/g_S)nh^2\le
\frac{99}{256}\hat\kappa nh^2$, since $L/g_S\le\kappa_S\le\hat\kappa$.

(iii) By induction on $j$, the center of stage $j$ satisfies
$\dist(c^{(j)},\mathcal S)^2\le4\hat\kappa nh_j^2$, so (ii) applies at stage
$j$. For $j=0$, $\dist(\ell,\mathcal S)^2\le ns^2=nh_0^2\le4\hat\kappa nh_0^2$.
If (ii) applies at stage $j$, then
$\dist(c^{(j+1)},\mathcal S)^2=\dist(y^{(j)},\mathcal S)^2
\le\frac{99}{256}\hat\kappa nh_j^2<2\hat\kappa nh_{j+1}^2$.

(iv) Let $\Delta$ be the common denominator of Section~\ref{sec:exact-data};
it also clears $s$ and $L$. By induction on $j$, every node of stage $j$ lies
in $\Gamma_j^{-1}\Z$ with $\Gamma_j=\Delta2^{j+\mu K_\theta}$: the center is
$\ell$ or a node of stage $j-1$; the box endpoints are $\ell_i$, $u_i$ or
$c_i\pm\varrho h_j$ with $h_j=s2^{-j}$; an unclipped node is $c_i\pm t_k$ with
$t_k=h_j\bigl[(2^\mu+1)^k-2^{\mu k}\bigr]/2^{\mu(k-1)}$ and $k<K_\theta$; and a
clipped node is a box endpoint. Hence, at every grid point, $\Delta\Gamma_j^2F$
is an integer (Section~\ref{sec:exact-data}), and so are $8\Delta\Gamma_j^2d_i$
and $4^{\mu+1}\Delta\Gamma_j^2\omega\norm{y-c}^2$. All table entries are
therefore integers over the common denominator $8\cdot4^\mu\Delta\Gamma_j^2$,
and their magnitudes are at most $2^{O(I)}$ times this denominator because
all points lie in $X$. Messages are sums of at most $\abs{\mathcal A}+2n$
such values, so all numbers have bit length $O(I+j+\mu K_\theta)$. The count
of table operations is that of Lemma~\ref{lem:dp} with at most $K_\theta$
nodes per coordinate.
\end{proof}

\begin{proof}[Proof of Theorem~\ref{thm:diagdiscovery}]
(a) DISC (Algorithm~\ref{alg:disc}) stops only when the hypotheses of
Lemma~\ref{lem:diagcert}(ii) hold, and it checks them in exact rational
arithmetic: the KKT conditions are sign tests, and positive semidefiniteness of $H+\Lambda(\bar x)$ is
decided by symmetric Gaussian elimination that rejects a negative pivot and
a zero pivot with a nonzero entry in its row. By
Lemma~\ref{lem:diagcert}(ii), $\bar x\in\mathcal S$, $\OPT=F(\bar x)$
and~\eqref{eq:diagset} holds. So an accepted $\bar x$ is a point of
$\mathcal S$ with $H+\Lambda(\bar x)\succeq0$, and then $F\in\mathfrak D$
by Lemma~\ref{lem:diagcert}(iii). Hence DISC does not stop if
$F\notin\mathfrak D$.

(b) Remark~\ref{rem:setgrowth} gives some $g_S>0$. Let $\hat\kappa\ge\kappa_S$
be a guess. By Lemma~\ref{lem:proximal}(iii) with $h_{J_{\hat\kappa}}^2
=s^24^{-J_{\hat\kappa}}\le\tau^2/(4\hat\kappa n)$,
$\dist(y,\mathcal S)^2\le\frac{99}{256}\cdot\frac{\tau^2}4<\frac{\tau^2}4$. By
Lemma~\ref{lem:snap}, REC$(y)$ (Algorithm~\ref{alg:rec}) does not fail, and
every point it can return, in particular $\bar x$, lies in $\mathcal S$; so $\bar x$ is a box KKT point.
Since $F\in\mathfrak D$, Lemma~\ref{lem:diagcert}(iii) gives
$H+\Lambda(\bar x)\succeq0$ for every $\bar x\in\mathcal S$, and the test
accepts, whichever optimal component contains $\bar x$. The guesses are
powers of two, so the first guess $\hat\kappa\ge\kappa_S\ge1$ satisfies
$\hat\kappa\le2\kappa_S$.

(c) By (b), at most $2+\log_2\kappa_S$ guesses run, each with
$\hat\kappa\le2\kappa_S$. For each,
$J_{\hat\kappa}+1=O(\log(\hat\kappa ns^2/\tau^2))+1=\poly(I)+O(\log\hat\kappa)$,
because $\tau=1/(4nR)$ and $\log R$ is polynomial in $I$. By
Lemma~\ref{lem:proximal}(i) and~(iv), a stage costs
$O(p(\abs{\mathcal A}+n+N)K_\theta^p)$ operations on numbers of bit length
$O(I+J_{\hat\kappa}+\mu K_\theta)$, where $\mu=O(\log(2\hat\kappa))$ and
$K_\theta\le48\sqrt{\hat\kappa}\lceil\log_2(n+2)\rceil$. By~\eqref{eq:logabsorb}
with $n$ in place of $n_P$, a bag table has at most
$K_\theta^p\le2(48p\sqrt{\hat\kappa})^p(n+2)\le2(68p\sqrt{\kappa_S})^p(n+2)$
entries, since $48\sqrt2<68$; and
$\mu K_\theta\le c\sqrt{\hat\kappa}\log(2\hat\kappa)\,I$ for an absolute
constant $c$. Every factor that depends only on $p$ and $\hat\kappa$ is
nondecreasing in $\hat\kappa$; its value at $\hat\kappa=2\kappa_S$ enters
$f'$, and for fixed $p$ these values are nondecreasing and polynomial in
$\kappa_S$. The remaining factors are polynomial in $I$ with an absolute
exponent. REC uses the original data, so its linear system has size
polynomial in $I$, and exact linear programming returns a vertex of
polynomial bit length in polynomial time
\cite[Theorem~6.4.12]{GrotschelLovaszSchrijver1988}. The KKT and
semidefiniteness tests are polynomial.
\end{proof}

\begin{remark}[A false guess]\label{rem:falseguess}
If $\hat\kappa<\kappa_S$, PROX$(\hat\kappa)$ (Algorithm~\ref{alg:prox}) may
return the same nonoptimal point at every stage. For
$F=x^2+y^2-3xy+\frac{63}{128}(x+y)$ on $[0,1]^2$, put $t=x+y$; then $F\ge t^2-\frac54t^2+\frac{63}{128}t$, a concave function of
$t\in[0,2]$, so $(1,1)$ is the unique minimizer with $\OPT=-\frac1{64}$. The
origin is a box KKT point with $F=0$, so every set-growth constant
satisfies $g_S\le(F(0)-\OPT)/\dist(0,\mathcal S)^2=\frac1{128}$, and
$\kappa_S\ge256$ since $L=2$. Exact computation shows that
PROX$(\hat\kappa)$ returns the origin at every stage $j\le33$ for
$\hat\kappa\in\{1,2\}$. For
$\hat\kappa=4$, the points $(0,0)$ and $(1,1)$ both minimize $Q_\omega$ at
stages $0$ and $1$, and PROX returns the origin at every stage $j\le33$ when
ties are broken towards the current center. These stages cover the budgets
of DISC: the least common denominator of Section~\ref{sec:exact-data} is
$\Delta=128$, so $R=2^{23}$, $\tau=2^{-26}$ and $J_{\hat\kappa}=28,28,29$ for
$\hat\kappa=1,2,4$. Lemma~\ref{lem:diagcert} rejects the origin:
$H+\Lambda=\bigl(\begin{smallmatrix}191/64&-3\\-3&191/64\end{smallmatrix}\bigr)$
has eigenvalue $-\frac1{64}$. At $(1,1)$ the matrix has diagonal
$\frac{193}{64}$ and is positive definite, so $F\in\mathfrak D$, and DISC
stops at a later guess by Theorem~\ref{thm:diagdiscovery}(b).
\end{remark}

\begin{proposition}[Discovery within $\mathfrak D$ is hard without the growth parameter]\label{prop:sshard}
For integers $a_0,a_1,\ldots,a_m$ with $A=\sum_{i=1}^m|a_i|\ge1$ and
$|a_0|\le A$, let $\xi_0=0$ and
\[
F(x,\xi)=\sum_{i=1}^m(\xi_i-\xi_{i-1}-a_ix_i)^2+(\xi_m-a_0)^2
+\sum_{i=1}^mx_i(1-x_i)
\]
on $x\in[0,1]^m$, $\xi\in[-A,A]^m$. The bags $\{\xi_1,x_1\}$ and
$\{\xi_{i-1},\xi_i,x_i\}$ $(2\le i\le m)$ form a path decomposition with bag
size at most $3$. If some $\hat x\in\{0,1\}^m$ has
$\sum_{i=1}^ma_i\hat x_i=a_0$, then $F\in\mathfrak D$, $\OPT=0$, and
$\mathcal S$ consists exactly of the points $(\hat x,\hat\xi)$ with $\hat x$
such a solution and $\hat\xi$ its partial sums. Otherwise $\OPT>0$.
Consequently, unless $\mathrm P=\mathrm{NP}$:
\begin{enumerate}[label=(\roman*)]
\item no polynomial-time algorithm returns an exact minimizer for every
instance of $\mathfrak D$ with bag size three, even though $\OPT=0$ is known
in advance for these instances;
\item there is no polynomial bound $\kappa_S\le\poly(I)$ that holds, for a
suitable set-growth constant, on every instance of this form in
$\mathfrak D$.
\end{enumerate}
\end{proposition}

The objective is the Subset Sum reduction of
\cite[Remark~2]{DelPiaKhajavirad2026} with unscaled partial sums; what the
proposition adds is that its feasible instances belong to $\mathfrak D$.

\begin{proof}[Proof of Proposition~\ref{prop:sshard}]
$F\ge0$ on the box, and $F=0$ exactly when all squares and all products
$x_i(1-x_i)$ vanish, that is, for binary $x$ with partial sums $\xi$ and
$\xi_m=a_0$; the partial sums lie in $[-A,A]$. If no solution exists, $F>0$
on the compact box. In the feasible case let $\hat z=(\hat x,\hat\xi)$ be
optimal. The sum of squares is a quadratic $\Phi$ with $\Phi(\hat z)=0$ and
$\nabla\Phi(\hat z)=0$, because each square is of an affine function
vanishing at $\hat z$; so $\Phi(z)=\frac12(z-\hat z)^{\T}H_{\mathrm{sq}}
(z-\hat z)$ with $H_{\mathrm{sq}}=\nabla^2\Phi\succeq0$. Also
$\sum_ix_i(1-x_i)=\frac12\sum_i2(x_i-0)(1-x_i)$. So
$F(z)-F(\hat z)=\frac12(z-\hat z)^{\T}H_{\mathrm{sq}}(z-\hat z)
+\frac12\sum_i2x_i(1-x_i)$. At $\hat z$ the gradient is $1-2\hat x_i=\pm1$ in
$x_i$ and $0$ in $\xi$, so $\hat z$ is a box KKT point with
$\lambda_i(\hat z)=2$ on $x$ and $0$ on $\xi$, and
$H+\Lambda(\hat z)=H_{\mathrm{sq}}\succeq0$. Thus $F\in\mathfrak D$.

Subset Sum with binary-encoded integers is NP-complete
\cite{GareyJohnson1979}, also when restricted to instances with $A\ge1$ and
$\abs{a_0}\le A$, because the other instances are decided trivially; and the
instance $F$ has size polynomial in that of the Subset Sum instance. (i) An
algorithm as in (i), stopped at its polynomial time bound and followed by a
check that its output $z$ is feasible with $F(z)=0$, would decide Subset Sum. (ii) If such a polynomial
bound existed, then by Theorem~\ref{thm:diagdiscovery}(c), whose bound is
nondecreasing and polynomial in $\kappa_S$ for fixed $p=3$, DISC would stop
within a polynomial time bound on every feasible instance; stopping DISC at
that bound and checking whether it has output $\OPT=0$ would decide Subset Sum.
\end{proof}

Unless $\mathrm P=\mathrm{NP}$, Corollary~\ref{cor:facecsp} has no
counterpart in $\mathfrak D$. Take $a_0=0$ in Proposition~\ref{prop:sshard}. The origin is then a minimizer, so
the instance belongs to $\mathfrak D$, and the origin is the only minimizer
exactly when no nonempty subset of $a_1,\ldots,a_m$ sums to zero. Deciding
whether such a subset exists is NP-complete: Subset Sum with positive items
$b_1,\ldots,b_k$ and target $T>0$ reduces to it by adding the item $-T$,
because a nonempty subset with sum zero must contain $-T$. Hence deciding
whether a given minimizer of an instance of $\mathfrak D$ with bag size three
is the only one is coNP-hard, even when the description~\eqref{eq:diagset} is
given.
